\section{Exercises}\label{sec:11-modules:07-exercises:1}

\textbf{Key points.} A module is an abelian group with a scalar action satisfying four axioms, equivalently a ring homomorphism into the ring of the endomorphisms of the group itself. A submodule is a membership predicate plus closure proofs, a linear map is a function preserving \texttt{+} and the scalar action, and the five/four axioms of a direct sum each split, via \texttt{congr 1}, into one independent fact per summand.

\begin{enumerate}
\def\labelenumi{\arabic{enumi}.}
\tightlist
\item
  Prove that the identity function is a linear map, for any \texttt{Mod : Module R Rg M}, construct \texttt{idLinearMap : LinearMap Rg Mod Mod} with \texttt{toFun := id}.
\item
  Prove that linear maps compose, given \texttt{f : LinearMap Rg ModM ModN} and \texttt{g : LinearMap Rg ModN ModP}, construct \texttt{LinearMap Rg ModM ModP} with \texttt{toFun := g.toFun ∘ f.toFun}. (This, together with Exercise 1 and associativity of function composition, is what makes \(R\)-modules and \(R\)-linear maps a category. This is the payoff that Chapter 1 promised.)
\item
  Verify that \texttt{intSmul} (defined above) satisfies the four axioms of \texttt{Module} against \texttt{intRing}, at least for \texttt{one\_\allowbreak{}smul} and \texttt{smul\_\allowbreak{}add}, by induction on the integer scalar, in the style of Chapter 5.
\item
  Define the submodule of multiples of a fixed integer \texttt{d : Int} (generalizing \texttt{evenSubmodule}, which is the case \texttt{d = 2}), and check that the three closure proofs generalize verbatim with \texttt{2} replaced by \texttt{d}.
\end{enumerate}

Solutions, \hyperref[sec:15-appendix-solutions:11-chapter-11:1]{Appendix, Chapter 11}.

